\subsection{The hybrid comparison: stopping and terminal information}
\label{subsec:r-hybrid-application}

The averaged block estimate combines with the common stopping
reduction.  We need a terminal information bound uniform over the
auxiliary row counts in that average.

\Needspace{6\baselineskip}
\begin{lemma}[Information in a successful terminal state]
\label{rhyb:terminal}
Let \(V\) be a padded final state whose unit-vector output has
uniform-prior angular success at least \(2/3\) at
\(0<\epsilon\le1/10\).  For every \(k_0\le k\le k_1\),
\begin{equation}\label{rhyb:terminal-information}
 J_k(V)\ge\frac{d-k-1}{14}\log\frac1{5\epsilon}-\log2.
\end{equation}
\end{lemma}

\begin{proof}
For every \(k_0\le k\le k_1\), the auxiliary matrix has
\(k\le d/2\) rows.  At \(\rho=2/5\),
\eqref{eq:r-residual-markov} therefore gives
\[
 \Pr\{R<2/5\}\le\frac{k/d}{1-4/25}\le\frac{25}{42}.
\]
The success probability is at least \(2/3\), and
\(2/3-25/42=1/14\).  Lemma~\ref{lem:r-residual-terminal}
now gives \eqref{rhyb:terminal-information}, because
\(\rho/(2\epsilon)=1/(5\epsilon)\).  Its radius condition
holds for \(\epsilon\le1/10\).
\end{proof}

\begin{theorem}[Precision lower bound from hybrid projections]
\label{rhyb:conclusion}
Let \(M(d)\) be a nonnegative integer sequence with \(M(d)=o(d^2)\), and
let \(0<\epsilon(d)\le1/10\).  For each \(d\), consider a learner in
Definition~\ref{def:r-streaming-model} with memory bound \(M(d)\),
accuracy \(\epsilon(d)\), deterministic horizon \(T(d)\ge0\), and
uniform prior on \(S^{d-1}\).  If
\[
 \mathbb P\{\angle(\widehat S,S)\le\epsilon(d)\}\ge\frac23,
\]
then
\[
 T(d)\ge c\,d\log\frac1{\epsilon(d)}
\]
for all sufficiently large \(d\).  The constant \(c>0\) is absolute;
the dimension threshold may depend on the memory sequence.
\end{theorem}

\begin{proof}
Apply Lemma~\ref{lem:r-seed-reduction} with \(p_0=2/3\) to the
original jointly measurable experiment.  Fix a seed satisfying
\eqref{eq:r-seed-factorization} with conditional uniform-prior success
at least \(2/3\), retain its conditional initial-state law, and then
apply Lemma~\ref{lem:r-borel-kernels}.  The fixed-seed joint experiment,
state bound, and horizon are preserved.

Put \(b=\log(1/\epsilon)\).  In dimensions with \(T>db\), the conclusion
already holds.  In the remaining dimensions, Lemma~\ref{lem:r-stopping-reduction}
permits padding with \(W=(T+2)2^M\) states and
\(B=\log W=o(d^2)\), so eventually \(B\le d^2\).
Break the padded computation into blocks of \(h_0\) observations, with
a shorter final block if necessary.  The initial padded state
\(\overline U_0\) is independent of \(S\), and the side matrix is
independent of the pair \((S,\overline U_0)\).  Thus
\(J_k(\overline U_0)=0\) for every \(k\), and its averaged potential is zero.
Proposition~\ref{rhyb:block} gives
\[
 \overline J(V_{\rm final})
 \le Cd\left\lceil\frac{T}{h_0}\right\rceil
 \le C'(T+d)\le C''T.
\]
The last inequality uses the independent bound \(T>d/4\) from
Lemma~\ref{lem:r-stopping-reduction}; thus it also covers a shorter
last block.

Lemma~\ref{rhyb:terminal} gives the opposite bound for every \(k\) in
the average.  Since \(k\le d/2\) and \(b\ge\log10\),
\[
 \log\frac1{5\epsilon}=b-\log5
 \ge\frac{\log2}{\log10}\,b.
\]
The right side of \eqref{rhyb:terminal-information} is consequently at
least \(c'db\) for all sufficiently large \(d\), uniformly over the
allowed accuracies.  Comparing the two bounds proves \(T\ge cdb\).
The dimensions with \(T>db\) were already covered.
\end{proof}

If a learner succeeds for every unit signal with probability at least
\(2/3\), it also satisfies the uniform-prior hypothesis, so the theorem
applies to that worst-signal guarantee.  The shared-seed case is covered
by the joint-independence reduction at the start of the proof, which
retains fresh randomized transition and output kernels.  The proof uses
exact labels and imposes no linearity or within-transition precision
restriction.
